\documentclass[12pt, letterpaper]{article}
\usepackage[mmddyyyy]{datetime}

\usepackage[utf8]{inputenc}
\usepackage{amsmath}
\usepackage{mathpazo}
\usepackage{mathtools}
\usepackage{eulervm}
\usepackage[skip=10pt plus1pt]{parskip}
\usepackage{textalpha}

\usepackage[margin=1in]{geometry}
\linespread{1.20}
\usepackage{enumerate}
\usepackage{enumitem}
\usepackage{listings}
\usepackage[T1]{fontenc}
\usepackage{amsfonts}
\usepackage{amssymb}
\usepackage{amsthm}
\usepackage{multirow,array}
\usepackage{graphicx}
\usepackage{physics}
\usepackage{bm}
\usepackage{tikz}
\usepackage[longnamesfirst]{natbib}
\usepackage{comment}
\usepackage{setspace}
\usepackage{hyperref}
\usepackage{accents}
\usepackage{graphicx}
\usepackage{float}
\usepackage{fancyvrb}
\usepackage{cancel}

\VerbatimFootnotes

\graphicspath{{Images/}}
\setlength{\parindent}{0pt}

\newtheorem{theorem}{Theorem}
\newtheorem{corollary}{Corollary}
\newtheorem{lemma}{Lemma}
\newtheorem{definition}{Definition}
\newtheorem{axiom}{Axiom}
\newcommand{\e}{\mathrm{e}}
\newcommand{\Var}{\mathrm{Var}}
\newcommand{\Ex}[1]{\mathbb{E}\left[#1\right]}
\renewcommand{\Pr}[1]{\mathbb{P}(#1)}
\newcommand{\plim}{\mathrm{plim}}
\newcommand{\pc}[1]{\frac{\dot{#1}(t)}{#1(t)}}
\renewcommand\qedsymbol{$\blacksquare$}
\newcommand{\ubar}[1]{\underaccent{\bar}{#1}}

\title{Homework 6 \\ Macroeconomic Theory I}
\author{Alejandro Flores-Brown}

\begin{document}

    \maketitle

    \section{Finishing Up Misallocation}

    \section{Alternative Sources for Variable Markups}

    \subsection{}

    The representative household solves

    \[ \max_{q_i} \int_0^1 p_i q_i di \text{ s.t. } 1 - \int_0^1 G\left(\frac{q_i}{Q}\right) di \geq 0    \]

    We use this problem to formulate the Lagrangian

    \[ \mathcal{L} = \int_0^1 p_i q_i di + \lambda \left[ 1 - \int_0^1 G\left(\frac{q_i}{Q}\right) di \right]  \]

    and take the first-order condition with respect to \(q_i\) to find

    \[ p_i = \lambda \frac{1}{Q} G'\left(\frac{q_i}{Q}\right)  \]

    \subsection{}

    Now we want to find an aggregate price index. We start this by multiplying both sides by \(q_i\) to get

    \[ p_i q_i = \lambda \frac{q_i}{Q} G'\left(\frac{q_i}{Q}\right)   \]

    and then integrate both sides over the measure of quantities to get 

    \[ \int_0^1 p_i q_i di = \lambda \int_0^1 \frac{q_i}{Q} G'\left(\frac{q_i}{Q}\right) di = \lambda   \]

    using the additional condition we are given at the beginning of the problem. It remains to solve the integral over \(p_i q_i\). However, because this is a representative household, we know that the sum of all prices and quantities purchased \emph{should} just be the aggregate prices and quantities! So we can write

    \[ \int_0^1 p_i q_i di = P Q = \lambda  \]

    which we obviously could rearrange into the desired form 

    \[ P = \int_0^1 p_i \frac{q_i}{Q} di    \]

    by dividing both sides by \(Q\). Now we plug in \(PQ\) for \(\lambda\) to get 

    \[ p_i = (P Q) \frac{1}{Q} G'\left(\frac{q_i}{Q}\right) \implies \frac{p_i}{P} = G'\left(\frac{q_i}{Q}\right)   \]

    as required.

    \subsection{}

    Now we consider the function \(G(x)\) to be

    \[ G(x) = 1 + \frac{1 - \e^{-\alpha(x - 1)}}{\alpha} = 1 + \frac{1}{\alpha} - \frac{1}{\alpha} \e^{-\alpha(x - 1)}   \]

    We will make our lives easier by immediately finding the derivative:

    \[ G'(x) = - \frac{1}{\alpha} \frac{d}{dx} \left(  \e^{-\alpha(x - 1)}  \right) = - \frac{1}{\alpha} \e^{-\alpha (x - 1)} (- \alpha) = \e^{-\alpha (x - 1)}   \]

    Now we derive the demand for good \(i\). This is given by \(q_i\), which we can find in the first-order condition we found in the previous part:

    \[ \frac{p_i}{P} = G'\left(\frac{q_i}{Q}\right) = \e^{-\alpha (\frac{q_i}{Q} - 1)} \implies \ln(\frac{p_i}{P}) = - \alpha \frac{q_i}{Q}     \]

    Rearranging this condition, we recover the demand function

    \[ q_i = Q \left( 1 + \frac{1}{\alpha} \ln(\frac{p_i}{P}) \right)      \]

    \subsection{}

    Now we check our first-order condition. Allowing \(p_i = P\) we recover

    \[ q_i = Q \left(1 + \frac{1}{\alpha} \ln(1) \right) = Q   \]

    like we want.

    \subsection{}

    We are given 

    \[ \epsilon_i \equiv \frac{\partial \ln(q_i)}{\partial \ln(p_i)} = - \frac{\partial q_i}{\partial p_i} \frac{p_i}{q_i}    \]

    Taking the condition for demand we found before, we can take the derivative to find

    \[ \frac{\partial q_i}{\partial p_i} = \frac{partial}{\partial p_i} \left(Q + \frac{Q}{\alpha } \ln(\frac{p_i}{P})\right) = \frac{Q}{\alpha} \left[ \ln(P) - \ln(p_i) \right] = - \frac{Q}{\alpha} p_i^{-1}     \]

    This then lets us recover the demand elasticity as

    \[ \epsilon_i = - \frac{\partial q_i}{\partial p_i} \frac{p_i}{q_i} = \frac{Q}{\alpha} p_i^{-1} \frac{p_i}{q_i} = \frac{1}{\alpha} \left(\frac{q_i}{Q}\right)^{-1}  \]



    \section{Two Country Atkeson-Burstein}

\end{document}